\documentclass[10pt,twocolumn]{article}
\usepackage[margin=0.85in]{geometry}
\IfFileExists{lmodern.sty}{\usepackage[T1]{fontenc}\usepackage{lmodern}\usepackage{microtype}}{}
\usepackage{amsmath}
\usepackage{booktabs}
\usepackage{xcolor}
\usepackage{hyperref}
\usepackage{fancyhdr}
\hypersetup{colorlinks=true,linkcolor=black,citecolor=black,urlcolor=black}

% Every result-dependent passage is marked with \pending{...}.
% Rule: a \pending box is replaced only with numbers produced by
% analysis/analyze.py from raw logs committed under experiments/.
\newcommand{\pending}[1]{\par\noindent\fbox{\parbox{\dimexpr\linewidth-2\fboxsep-2\fboxrule}{\small\textbf{PENDING MEASURED DATA.} #1}}\par}

\pagestyle{fancy}
\fancyhf{}
\fancyhead[C]{\small DRAFT --- methods complete, results pending measured runs. Not for citation.}
\fancyfoot[C]{\thepage}
\renewcommand{\headrulewidth}{0pt}

\title{What Does a Streaming Client See When an LLM Server Crashes?\\
\large A Fault-Injection Study of Client-Visible Token Anomalies in vLLM and SGLang}
\author{Rajshekar Medipally\\
\small Independent researcher, Richmond, VA\\
\small \url{https://github.com/rmedipallycic/llm-inference-fault-tolerance}}
\date{Draft, October 2026}

% Table rows come from paper/tables/*.tex, written by analysis/make_tables.py
% from the committed raw logs. Until those files exist, placeholders show.
% (\@@input is the expandable primitive, so it is safe inside tabular.)
\makeatletter
\IfFileExists{tables/anomalies.tex}{\def\AnomRows{\@@input tables/anomalies.tex }}{\def\AnomRows{%
vLLM   & naive    & 0   & -- & -- & -- \\
vLLM   & naive    & 0.8 & -- & -- & -- \\
vLLM   & recomp.  & 0   & -- & -- & -- \\
vLLM   & recomp.  & 0.8 & -- & -- & -- \\
SGLang & naive    & 0   & -- & -- & -- \\
SGLang & naive    & 0.8 & -- & -- & -- \\
SGLang & recomp.  & 0   & -- & -- & -- \\
SGLang & recomp.  & 0.8 & -- & -- & -- \\
}}
\IfFileExists{tables/latency.tex}{\def\LatRows{\@@input tables/latency.tex }}{\def\LatRows{%
vLLM   & -- & -- \\
SGLang & -- & -- \\
}}
\IfFileExists{tables/env.tex}{\def\EnvRows{\@@input tables/env.tex }}{\def\EnvRows{}}
\makeatother

\begin{document}
\maketitle
\thispagestyle{fancy}

\begin{abstract}
LLM inference servers stream generated tokens to clients over requests that
last seconds to minutes. When a serving process fails mid-generation, the
client is left holding a partial response, and the recovery it performs
determines what it ultimately shows to a user or passes to a downstream agent.
Prior work on fault-tolerant LLM serving measures server-side recovery time.
This paper instead measures what the client receives. We define three
client-visible anomalies (duplicated, missing, and divergent text) relative to
a failure-free reference, build an open-source fault-injection harness that
SIGKILLs an unmodified OpenAI-compatible server at a chosen token position and
records every streamed chunk, and apply it to vLLM and SGLang serving
Llama~3.1~8B~Instruct on a single GPU. We compare the two recovery strategies
available to a client of an unmodified engine: naive retry, which resubmits
the original prompt, and recompute, which resubmits the prompt plus the text
already delivered.
\pending{One to two sentences stating the measured anomaly rates and recovery
latencies for each engine, strategy, and temperature.}
All raw logs and the single analysis script are public.
\end{abstract}

\section{Introduction}

LLM serving has become stateful streaming infrastructure. A generation request
holds a growing key-value (KV) cache on an accelerator and emits tokens to the
client incrementally. Failure handling, however, is still largely stateless:
if the serving process dies, the client or a load balancer resubmits the
request and generation starts again.

That recovery has a cost that server-side metrics miss. A client that has
already displayed 128 tokens may receive those 128 tokens a second time, may
receive a continuation that disagrees with what it already showed, or may lose
text at the boundary. For a chat interface this is a visible glitch. For an
agent that has already acted on a partial tool call or a partial plan, it is a
correctness problem. Stream processing systems name this property
\emph{delivery semantics} and engineer exactly-once output deliberately
\cite{elnozahy2002survey,carbone2015lightweight,armbrust2018structured}. LLM
serving has no equivalent definition, and to our knowledge no measurement of
what clients actually receive.

This paper makes three contributions:
\begin{enumerate}
  \item \textbf{Definitions.} Token-level delivery semantics for streamed
    generation, and three measurable client-visible anomalies: duplication,
    loss, and divergence (Section~\ref{sec:model}).
  \item \textbf{A fault-injection harness} for any OpenAI-compatible server,
    which kills the full server process tree at a chosen chunk, restarts it,
    applies a client recovery strategy, and records every chunk with
    timestamps, the code commit, package versions, and GPU inventory
    (Section~\ref{sec:harness}).
  \item \textbf{A measurement study} of naive retry and recompute on vLLM and
    SGLang at temperature 0 and 0.8, at three kill positions
    (Sections~\ref{sec:method}--\ref{sec:results}).
\end{enumerate}

The study deliberately uses unmodified engines and client-side recovery only.
It establishes the baseline that a server-side exactly-once protocol, the next
stage of this project, would have to improve on.

\section{Background}

\paragraph{Autoregressive serving.}
A request runs a prefill phase, which processes the prompt in parallel and
builds the KV cache, followed by a decode phase that produces one token per
step and streams it to the client. Engines such as vLLM \cite{kwon2023pagedattention}
and SGLang \cite{zheng2024sglang} batch requests at iteration granularity
\cite{yu2022orca} and manage the KV cache in fixed-size blocks.

\paragraph{The state lost on failure.}
The KV cache is a deterministic function of the weights and the token
sequence, so it can always be rebuilt by prefill. What cannot be rebuilt from
the prompt alone is the \emph{sampling path}: at temperature above zero, the
continuation depends on random draws whose state lives in the dead process.

\paragraph{Rollback-recovery vocabulary.}
Classic rollback-recovery distinguishes checkpointing from logging and defines
\emph{output commit}: a message to the outside world must not be released
until the state that produced it is recoverable \cite{elnozahy2002survey}.
Mapped onto LLM serving, the token sequence is a replayable input log, the KV
cache is derived operator state, and the streaming client is the external
sink where delivery semantics are decided.

\section{Delivery model and anomalies}
\label{sec:model}

Let $G$ be the failure-free reference output for a request, recorded on the
same server configuration. A failure occurs after the client has received
delivered text $D$. After recovery the client receives recovery text $R$ and
its final view is $V = D \cdot R$ (concatenation), which is what a typical
client displays.

We measure:
\begin{itemize}
  \item \textbf{Duplication}: the length of the longest common prefix of $R$
    and $D$, reported as the fraction $\textit{dup} = |\mathrm{lcp}(R, D)| / |D|$.
    $\textit{dup}=1$ means the client received all of $D$ a second time.
  \item \textbf{Divergence}: whether the client's text disagrees with $G$
    after the duplicated prefix is removed, i.e.\ whether
    $V' = D \cdot R[\textit{dup}{:}]$ and $G$ differ at some position before
    either ends. Removing the duplicate first keeps the two anomalies
    separate: a retry that re-sends identical text is duplication, not
    divergence. The first differing index is recorded and compared with
    $|D|$, the kill boundary.
  \item \textbf{Exact match}: $V = G$, i.e.\ the failure was invisible in the
    final text.
  \item \textbf{Recovery latency}: wall time from the kill to the first chunk
    delivered after recovery, including server restart.
\end{itemize}

Under this model, \emph{at-least-once} delivery shows as $\textit{dup}>0$,
\emph{at-most-once} as $V$ being a strict prefix of $G$, and
\emph{exactly-once} as $V=G$. Comparison is on text, not token ids, because
the OpenAI-compatible streaming API does not guarantee one token per chunk.

A comparison against $G$ is only meaningful if the server is reproducible
without failures. The harness therefore records two failure-free references per
run and flags the run when they differ (Section~\ref{sec:harness}).

\paragraph{Client recovery strategies.}
\textbf{Naive retry} resubmits the original prompt with the original
\texttt{max\_tokens}; many client libraries' retry wrappers do this.
\textbf{Recompute} resubmits the prompt concatenated with $D$ and requests the
remaining token budget, so the server rebuilds the KV cache by prefill and
continues. Neither requires engine changes.

\section{Fault-injection harness}
\label{sec:harness}

The harness (\texttt{harness/}, about 340 lines of Python) has three parts.

\paragraph{Server lifecycle.} The server command is started in its own process
group and declared ready when \texttt{GET /v1/models} returns 200. A fault is a
\texttt{SIGKILL} of the whole process group, since vLLM and SGLang spawn worker
processes that otherwise keep holding GPU memory and the port. On Windows,
used only for laptop smoke tests, the equivalent is \texttt{taskkill /F /T}.

\paragraph{Streaming client.} Requests go to \texttt{/v1/completions} with
\texttt{stream=true}. Every server-sent event with non-empty text is recorded
as a chunk with wall-clock and monotonic timestamps. The kill fires from the
chunk callback when the configured chunk count is reached, and the client
keeps reading, so chunks the server had already sent before dying are recorded
as in-flight rather than lost.

\paragraph{Trial protocol.} Each run (one engine, strategy, temperature, and
kill position) starts the server, records two failure-free references $G$ and
$G'$, then for each trial: streams the request, kills at the kill point,
restarts the server and waits for readiness, issues the recovery request,
and writes a JSON record containing every chunk, the kill time, the restart
time, and the client view $V$.

\paragraph{Provenance.} Each run records the git commit and whether the
working tree was dirty, the full command line, installed versions of
\texttt{vllm}, \texttt{sglang}, \texttt{torch}, and \texttt{transformers}, the
GPU model, memory, and driver from \texttt{nvidia-smi}, and the full server
log. Measured runs are written to \texttt{experiments/} and committed; harness
smoke tests are written to a separate, git-ignored directory and are never
reported. All statistics in this paper are produced by one script,
\texttt{analysis/analyze.py}, from those committed files.

\paragraph{Validation.} The harness was exercised end to end on a laptop CPU
with llama.cpp and a 0.5B model before any GPU run. That check, documented in
the repository as a mechanics check rather than a result, led to two changes:
a duplication metric based on shared prefix fraction rather than any shared
character, and the reference-reproducibility check, after two clean seeded
runs on that server disagreed at temperature 0.8.

\section{Methodology}
\label{sec:method}

\paragraph{Setup.} One cloud machine with a single NVIDIA GPU of at least
24\,GB, Llama~3.1~8B~Instruct in bf16, vLLM and SGLang each in a separate
Python environment with default engine flags except a 4{,}096-token
context cap (Llama~3.1's 128k default does not fit a 24\,GB KV cache; every
request here is far shorter). The GPU model, driver, and
package versions are recorded by the harness in every run and reported as
recorded\IfFileExists{tables/env.tex}{ (Table~\ref{tab:env})}{}.

\IfFileExists{tables/env.tex}{%
\begin{table}[h]
\centering\small
\caption{Software and hardware, from \texttt{run.json}.}
\label{tab:env}
\begin{tabular}{@{}lllp{0.38\linewidth}@{}}
\toprule
Engine & Version & torch & GPU \\
\midrule
\EnvRows
\bottomrule
\end{tabular}
\end{table}
\IfFileExists{tables/provenance.tex}{\noindent\small\input{tables/provenance.tex}\par}{}%
}{\pending{Engine, torch, CUDA and driver versions and model revision hash,
generated from \texttt{run.json} by \texttt{analysis/make\_tables.py}.}}

\paragraph{Grid.} For each engine: strategy $\in$ \{naive retry, recompute\}
$\times$ temperature $\in \{0, 0.8\}$ $\times$ kill point $\in$ \{25, 128, 230\}
chunks (about 10\%, 50\%, and 90\% of a 256-token budget), with 10 trials
per cell and a fixed per-request seed (1234). That is 12 runs and 120 trials
per engine, 240 in total, produced by \texttt{scripts/run\_matrix.sh}.
Requests are issued one at a time, so there is no batching interference.

\paragraph{Hypotheses.} Stated in advance so that the data can refute them:
\begin{description}
  \item[H1] At temperature 0, naive retry yields $\textit{dup}=1$ in every
    trial and recompute yields $V=G$ in every trial.
  \item[H2] At temperature 0.8 with a fixed seed, both engines reproduce $G$
    across clean runs; if so, naive retry still yields full duplication and
    recompute yields divergence from $G$ after the kill point, because the
    sampler restarts from a different random state.
  \item[H3] Recovery latency is dominated by server restart, not by prefill
    of the delivered text.
\end{description}

\section{Results}
\label{sec:results}

\pending{\textbf{Reproducibility first.} For each engine and temperature,
report whether the two failure-free references matched (\texttt{reference
repeat identical}). Any cell where they did not is marked uninterpretable for
divergence and excluded from Table~\ref{tab:anomalies}, with the count stated.}

\begin{table}[t]
\centering\small
\caption{Client-visible anomalies per cell (10 trials each), generated from
the raw logs by \texttt{analysis/make\_tables.py}. Diverged$'$ is counted
after removing duplicated text (Section~\ref{sec:model}).}
\label{tab:anomalies}
\setlength{\tabcolsep}{3.5pt}
\begin{tabular}{@{}llrccc@{}}
\toprule
Engine & Strategy & $T$ & $V{=}G$ & Full dup & Diverged$'$ \\
\midrule
\AnomRows
\bottomrule
\end{tabular}
\end{table}

\begin{table}[t]
\centering\small
\caption{Restart and recovery latency, median (min--max), seconds.}
\label{tab:latency}
\begin{tabular}{@{}lcc@{}}
\toprule
Engine & Server restart & Kill $\to$ first recovered token \\
\midrule
\LatRows
\bottomrule
\end{tabular}
\end{table}

\pending{\textbf{H1, H2 outcomes} by kill position, with the raw counts.
State plainly where a hypothesis was refuted.}

\pending{\textbf{H3}: restart time versus recovery latency, and whether
recompute is measurably slower than naive retry to first token at the 90\%
kill point.}

\section{Discussion}

\pending{What the measured anomalies mean for chat clients and for agents
that act on partial output; which strategy a client should use for each
temperature; whether the difference between engines matters.}

\paragraph{Why client-side recovery is not enough.} Even in the best case for
recompute, the client cannot restore the sampler's random state, so above
temperature 0 it can at most hide duplication, not guarantee the continuation
the original request would have produced. Exactly-once delivery therefore
requires state the client does not hold: a sequence-numbered token log and the
sampler state, committed before tokens are released. That protocol is the
next stage of this project.

\section{Related work}

\textbf{Fault-tolerant LLM serving.} D\'ej\`aVu \cite{strati2024dejavu}
streams and replicates KV cache to recover from failures and reports
server-side recovery time; it does not characterize what streaming clients
receive. SpotServe \cite{miao2024spotserve} migrates in-flight context when
preemptible instances are reclaimed with warning, and Llumnix
\cite{sun2024llumnix} live-migrates requests for load balancing. Both handle
planned moves; this study targets unplanned fail-stop crashes.
\textbf{Serving engines.} Orca \cite{yu2022orca} introduced iteration-level
scheduling, vLLM \cite{kwon2023pagedattention} block-managed KV memory, and
SGLang \cite{zheng2024sglang} prefix sharing via RadixAttention.
\textbf{Delivery semantics.} Chandy and Lamport \cite{chandy1985snapshots}
define consistent snapshots, Elnozahy et al.\ \cite{elnozahy2002survey}
output commit, and Flink \cite{carbone2015lightweight} and Spark Structured
Streaming \cite{armbrust2018structured} build exactly-once sinks on them.
Optimal checkpoint intervals \cite{young1974first,daly2006higher} motivate
later stages of this project.

\section{Threats to validity}

\begin{itemize}
  \item \textbf{One model, one GPU, one prompt.} Results are scoped to
    Llama~3.1~8B on one single-GPU machine with a single explanatory prompt. Tensor-parallel
    serving and other prompts may differ.
  \item \textbf{Re-tokenization in recompute.} Recompute sends $D$ as text, and
    the engine re-tokenizes it. If the boundary tokenizes differently than the
    original output tokens, the continuation can diverge at temperature 0 for
    a reason unrelated to sampling. Any such divergence is reported with its
    position relative to the kill point.
  \item \textbf{Fail-stop only.} \texttt{SIGKILL} is a clean crash; GPU hangs,
    NCCL stalls, and silent corruption are not modeled.
  \item \textbf{Latency includes restart.} With a single instance, recovery
    latency includes model reload; failover to a warm standby would be faster.
  \item \textbf{Chunks are approximate tokens.} Kill points are set in chunks,
    which are usually but not always one token.
\end{itemize}

\section{Conclusion and future work}

\pending{Two-sentence summary of measured findings.}

Next stages: (1) a gateway that assigns sequence numbers, logs tokens and
sampler state before release, and deduplicates on recovery, evaluated on the
same harness for zero anomalies and per-token overhead; (2) measured
recompute-versus-checkpoint cost curves for the KV cache; (3) an adaptive
per-request checkpoint policy.

\paragraph{Artifact.} Code, raw logs, and the analysis script are at
\url{https://github.com/rmedipallycic/llm-inference-fault-tolerance}
(Apache-2.0). Every number above traces to a committed log file.

\bibliographystyle{plain}
\bibliography{refs}
\end{document}
